% ================================================================= MỞ RỘNG SANG SILESIA
\subsection{Mở rộng sang 22 sản phụ và sang thai kỳ}
\label{sec:silesia}

Điểm yếu lớn nhất của bản đề cương trước là chỉ có 5 sản phụ, tất cả đang chuyển dạ. Nhóm đã tải bộ
Silesia (Matonia và cộng sự, \textit{Scientific Data} 7:200, 2020; figshare DOI 10.6084/m9.figshare.c.4740794,
giấy phép CC0), gồm \textbf{10 sản phụ thai kỳ 32--42 tuần} (B1, mỗi người 20 phút) và \textbf{12 sản phụ chuyển
dạ} (B2, mỗi người 5 phút), cùng bệnh viện và cùng hệ ghi KOMPOREL với dữ liệu huấn luyện. Toàn bộ do
\texttt{model/download\_silesia.py} tải với mã băm SHA-256 ghi lại.

\subsubsection{Kiểm tra rò rỉ trước khi đánh giá}

Bài báo gốc tự khai rằng 5 bản ghi PhysioNet được lấy từ bộ chuyển dạ. Nhóm không tin lời khai mà đo lại
bằng tín hiệu: đưa cả hai về 250\,Hz, trượt cửa sổ 60 giây của từng bản ghi PhysioNet trên toàn bộ từng
bản ghi B2, thử mọi tổ hợp 4$\times$4 đạo trình và mọi độ lệch $\pm$5 phút.

\begin{table}[H]\centering\small
\caption{Kết quả kiểm tra rò rỉ. NCC là tương quan chéo chuẩn hoá cực đại. Cột ``hạng nhì'' là NCC với
bản ghi B2 gần nhất còn lại, để thấy khoảng cách.}
\label{tab:roriSilesia}
\begin{tabular}{@{}C{1.6cm} C{1.6cm} C{1.6cm} C{2.0cm} C{1.8cm} C{2.2cm} C{2.2cm}@{}}
\toprule
\textbf{PhysioNet} & \textbf{Silesia} & \textbf{NCC} & \textbf{Lệch} & \textbf{Hạng nhì} &
\textbf{Nhãn khớp} & \textbf{Số nhịp} \\
\midrule
r01 & B2\_01 & 0,989 & $-56$\,ms & 0,135 & 100,0\% & 644 / 644 \\
r10 & B2\_02 & 0,993 & $-40$\,ms & 0,147 & 99,8\% & 637 / 637 \\
r04 & B2\_07 & 0,994 & $-48$\,ms & 0,178 & 99,8\% & 632 / 632 \\
r07 & B2\_10 & 0,994 & $-48$\,ms & 0,141 & 99,8\% & 627 / 627 \\
r08 & B2\_11 & 0,988 & $-56$\,ms & 0,128 & 98,8\% & 651 / 646 \\
\bottomrule
\end{tabular}
\end{table}

Kết luận: 5/12 bản ghi B2 \textbf{chính là} 5 bản ghi PhysioNet. Hệ quả thứ nhất: số sản phụ độc lập
là \textbf{22}, không phải 27. Hệ quả thứ hai: với 5 bản ghi trùng, nhóm dùng checkpoint fold tương ứng
(chưa từng thấy bản ghi đó); 7 bản ghi B2 còn lại và 10 bản ghi B1 dùng checkpoint production, zero-shot.

\subsubsection{Kết quả}

\begin{table}[H]\centering\small
\caption{Cùng một mô hình, cùng giao thức $\pm$50\,ms, trên năm cấu hình dữ liệu. ``$\geq$90'' và ``$<$50''
là số bản ghi có F1 (kênh PSD) trong khoảng đó.}
\label{tab:silesia_tonghop}
\begin{tabular}{@{}L{5.0cm} C{0.8cm} C{1.2cm} C{2.6cm} C{2.6cm} C{0.8cm} C{0.8cm}@{}}
\toprule
\textbf{Bộ dữ liệu} & $n$ & \textbf{Phút} & \textbf{F1, kênh PSD} & \textbf{F1, 4 kênh} & $\geq$90 & $<$50 \\
\midrule
ADFECGDB PhysioNet, tách bản ghi & 5 & 25 & $99{,}21 \pm 1{,}50$ & $97{,}45 \pm 2{,}96$ & 5 & 0 \\
Silesia B2 chuyển dạ, 12 bản ghi & 12 & 60 & $97{,}17 \pm 7{,}45$ & $95{,}26 \pm 9{,}59$ & 11 & 0 \\
\quad chỉ 7 bản ghi hoàn toàn mới, zero-shot & 7 & 35 & $95{,}87 \pm 9{,}75$ & $93{,}73 \pm 12{,}48$ & 6 & 0 \\
\textbf{Silesia B1 thai kỳ 32--42 tuần, zero-shot} & 10 & 199,6 & $\mathbf{93{,}30 \pm 13{,}46}$ &
$\mathbf{91{,}31 \pm 9{,}99}$ & 8 & 0 \\
CinC 2013 set-a, zero-shot & 10 & 10 & $59{,}15 \pm 37{,}17$ & $61{,}96 \pm 34{,}09$ & 4 & 5 \\
\bottomrule
\end{tabular}
\end{table}

\begin{table}[H]\centering\small
\caption{Silesia B1 --- thai kỳ, từng bản ghi. Mô hình chưa từng thấy bất kỳ dữ liệu thai kỳ nào.
Cột ``oracle'' là kênh tốt nhất theo nhãn, chỉ để tham chiếu.}
\label{tab:silesiaB1}
\begin{tabular}{@{}C{1.3cm} C{1.3cm} C{1.5cm} C{1.4cm} C{1.4cm} C{1.6cm} C{1.6cm} C{1.5cm}@{}}
\toprule
\textbf{Bản ghi} & \textbf{Kênh PSD} & \textbf{F1 PSD} & \textbf{Se} & \textbf{PPV} & \textbf{Jitter} &
\textbf{F1 4 kênh} & \textbf{Oracle} \\
\midrule
B1\_01 & A3 & 97,70 & 96,70 & 98,72 & 10,3\,ms & 87,79 & 97,70 \\
B1\_02 & A3 & 99,59 & 99,50 & 99,68 & 9,1\,ms & 97,97 & 99,59 \\
B1\_03 & A3 & 99,86 & 99,88 & 99,84 & 2,5\,ms & 99,17 & 99,86 \\
B1\_04 & A3 & 99,80 & 99,64 & 99,96 & 3,8\,ms & 98,59 & 99,96 \\
B1\_05 & A3 & 99,78 & 99,75 & 99,82 & 3,6\,ms & 99,47 & 99,78 \\
B1\_06 & A4 & 80,90 & 78,05 & 83,95 & 12,5\,ms & 92,43 & 98,28 \\
B1\_07 & A4 & \xau{58,72} & 53,21 & 65,49 & 15,3\,ms & 67,96 & 78,52 \\
B1\_08 & A1 & 99,52 & 99,28 & 99,76 & 7,3\,ms & 98,27 & 99,52 \\
B1\_09 & A3 & 99,04 & 98,92 & 99,16 & 4,6\,ms & 86,66 & 99,04 \\
B1\_10 & A3 & 98,12 & 98,46 & 97,78 & 13,2\,ms & 84,76 & 98,12 \\
\midrule
\textbf{Macro} & & $\mathbf{93{,}30 \pm 13{,}46}$ & 92,34 & 94,42 & 8,2\,ms & $91{,}31 \pm 9{,}99$ & $97{,}04 \pm 6{,}56$ \\
\bottomrule
\end{tabular}
\end{table}

Gộp toàn cục trên B1: 26.085 nhịp đúng, 1.451 báo động giả, 2.320 nhịp sót trên 28.405 nhịp có nhãn.

\subsubsection{Đọc kết quả}

\begin{enumerate}[leftmargin=1.6em,itemsep=4pt]
\item \textbf{Mô hình tổng quát được sang thai kỳ.} Huấn luyện chỉ trên 5 ca chuyển dạ, chưa từng thấy dữ
liệu thai kỳ, mà đạt 93,30 trên 10 sản phụ tuần 32--42. Tám trên mười bản ghi trên 90, \textbf{không bản
ghi nào dưới 50}. Phân bố này \textit{không} lưỡng cực như CinC 2013.
\item \textbf{Sụp đổ trên CinC 2013 không phải do giai đoạn thai.} Chuyển từ chuyển dạ sang thai kỳ trong
cùng hệ ghi mất khoảng 4 điểm ($97 \to 93$). Chuyển sang hệ ghi khác mất 38 điểm ($97 \to 59$). Biến gây
sụp đổ là \textbf{thiết bị và bố trí điện cực}, không phải tuổi thai.
\item \textbf{Quy tắc chọn kênh PSD yếu đi khi tín hiệu thai yếu.} Hai bản ghi kém nhất (B1\_06, B1\_07) đều
do quy tắc PSD chọn đúng kênh xấu nhất; với kênh tốt nhất chúng đạt 98,28 và 78,52. B1\_07 là bản ghi có
tín hiệu thai yếu nhất bộ theo chính tác giả. Đây là bằng chứng thứ hai (sau CinC) rằng heuristic phổ
cần được thay bằng cổng học được.
\item \textbf{Bản ghi B2 yếu duy nhất có lý do ghi trong tài liệu gốc.} B2\_03 đạt 73,81; tác giả ghi
``mất tín hiệu 27,5\%'' cho nhịp tim tham chiếu của chính bản ghi này.
\end{enumerate}

\begin{canhbao}[title={Nhãn B1 không phải chuẩn vàng --- phải nói rõ khi công bố}]
Thai kỳ không thể gắn điện cực da đầu. Nhãn nhịp thai của B1 do tác giả tạo bằng cách khử ECG mẹ, dò tự
động, rồi chuyên gia sửa tay và gắn cờ tin cậy (0,46\% nhịp cờ 0; bỏ chúng đi kết quả không đổi). Vì
vậy F1 trên B1 đo \textit{mức đồng thuận với pipeline của tác giả}, không phải sự thật sinh lý. Nhóm còn
đo được nhãn B1 lệch hệ thống 8--12\,ms sau đỉnh R bụng ở 5/10 bản ghi; lệch này nằm trong dung sai 50\,ms
nên không ảnh hưởng F1, nhưng jitter trên B1 (8,2\,ms) không so sánh được với B2 (2,9\,ms).

Silesia cùng bệnh viện, cùng hệ ghi, cùng bố trí điện cực với dữ liệu huấn luyện. Kết quả này chứng minh
tổng quát sang \textbf{giai đoạn thai kỳ}, chưa chứng minh tổng quát sang \textbf{hệ thống ghi khác}.
Các con số trong mục này là của mô hình 5 ca; mô hình huấn luyện lại trên 22 ca ở \S\ref{sec:model22}.
\end{canhbao}
